\subsection{Research Questions}

Building on the Moltbook Illusion and the gap between laboratory and deployment norm emergence, we address three research questions:

\begin{description}
\item[RQ1:] \textbf{What environmental factors determine whether norms emerge in LLM multi-agent populations?} We hypothesize that communication structure, not model capability, is the dominant determinant. We test this through a fully crossed $2\times2\times2\times2$ factorial experiment varying four environmental dimensions.

\item[RQ2:] \textbf{What are the distinct trajectories and failure modes of norm emergence in LLM populations?} Beyond whether norms emerge, we characterize \emph{how} they emerge—their dynamics, rates, and conditions under which they fail to crystallize. We identify four distinct emergence trajectories and model their boundary conditions.

\item[RQ3:] \textbf{How do agent-level properties (theory-of-mind reasoning, persona consistency, prompt framing) moderate norm emergence?} Even if environment dominates, agent properties may modulate the speed and stability of norm emergence. We measure the interaction between environmental configuration and agent-level factors.
\end{description}

These questions operationalize the broader HCI challenge: understanding when, how, and why social norms emerge in multi-agent AI systems, and what platform designers can do to shape emergence intentionally.